\DocumentMetadata{
  pdfversion=2.0,
  pdfstandard=ua-2,
  lang=en-US,
  tagging=on,
  tagging-setup={math/setup=mathml-SE,tabsorder=structure}
}
\documentclass[11pt,letterpaper]{article}
\usepackage[margin=0.68in,includefoot]{geometry}
\usepackage{newcomputermodern}
\usepackage{amsmath,amscd,mathtools}
\usepackage{tikz,adjustbox}
\providecommand{\square}{\mdlgwhtsquare}
\usepackage[hidelinks]{hyperref}
\hypersetup{
  pdftitle={Topology First-Year Exam, January 2025, Part 1},
  pdfauthor={University of Florida Department of Mathematics},
  pdfsubject={Graduate mathematics examination},
  pdfdisplaydoctitle=true,
  bookmarksopen=true
}
\setcounter{secnumdepth}{0}
\setlength{\parindent}{0pt}
\setlength{\parskip}{0.28em}
\setlength{\emergencystretch}{3em}
\setlength{\leftmargini}{2.15em}
\setlength{\leftmarginii}{2.65em}
\setlength{\leftmarginiii}{3em}
\renewcommand{\labelenumi}{\textbf{\theenumi.}}
\pagestyle{empty}
\raggedbottom
\allowdisplaybreaks
\newenvironment{examproblems}[1]{%
  \begin{enumerate}%
  \setcounter{enumi}{#1}%
  \setlength{\topsep}{0.35em}%
  \setlength{\partopsep}{0pt}%
  \setlength{\itemsep}{0.48em}%
  \setlength{\parsep}{0.18em}%
}{\end{enumerate}}
\newenvironment{examsubparts}{%
  \begin{enumerate}%
  \setlength{\topsep}{0.18em}%
  \setlength{\partopsep}{0pt}%
  \setlength{\itemsep}{0.16em}%
  \setlength{\parsep}{0.1em}%
}{\end{enumerate}}
\newenvironment{examinstructions}{%
  \par\begingroup\itshape
}{%
  \par\endgroup\vspace{0.18em}
}
\makeatletter
\renewcommand\section{\@startsection{section}{1}{0pt}%
  {-1.25ex plus -.25ex minus -.1ex}%
  {0.55ex plus .1ex}%
  {\normalfont\large\bfseries}}
\renewcommand{\@maketitle}{%
  \newpage\null\vskip -1em%
  {\footnotesize\bfseries
    \href{https://www.ufl.edu/}{University of Florida}, \href{https://math.ufl.edu/}{Department of Mathematics}\par}%
  \vskip -0.5em%
  \tagstructbegin{tag=Title}%
    {\LARGE\bfseries\href{https://gma.math.ufl.edu/exams/topology-first-year/}{Topology First-Year Exam}, January 2025, Part 1\par}%
  \tagstructend%
  \vskip 0.25em%
  \tagmcbegin{artifact}%
    \hrule height 1pt%
  \tagmcend%
  \vskip 1em%
}
\makeatother
\title{Topology First-Year Exam, January 2025, Part 1}
\author{}
\date{}
\begin{document}
\maketitle
\thispagestyle{empty}
\begin{examinstructions}
Answer the following problems and show all your work. Support all statements to the best of your ability.
\end{examinstructions}
\begin{examproblems}{0}
\item
Show that for every set~\(X\), there is no function \(f:X\to 2^X\) such that~\(f\) is onto. Note that \(2^X\) is the power set on~\(X\) and is also denoted by \(\mathcal{P}(X)\).
\item
Let \(\mathbb{R}^\infty\) denote the subset of \(\mathbb{R}^\omega\) consisting of all sequences that are “eventually zero,” that is, all sequences \((x_1,x_2,\ldots)\) such that \(x_i\ne 0\) for only finitely many values of~\(i\). What is the closure of \(\mathbb{R}^\infty\) in the box and product topologies? Justify your answer.
\item
Let \(p:X\to Y\) be continuous.
\begin{examsubparts}
\item[\textnormal{(a)}]
Define what it means for~\(p\) to be a quotient map.
\item[\textnormal{(b)}]
Show that if there is a continuous map \(f:Y\to X\) such that \(p\circ f\) equals the identity map of~\(Y\), then~\(p\) is a quotient map.
\end{examsubparts}
\item
This problem has two parts.
\begin{examsubparts}
\item[\textnormal{(a)}]
Give the topology of the one-point compactification of a locally compact Hausdorff topological space~\(X\).
\item[\textnormal{(b)}]
Show that the one-point compactification of \(\mathbb{Z}_+\) is homeomorphic with the subspace \(\{0\}\cup\{\frac{1}{n}\mid n\in\mathbb{Z}_+\}\) of~\(\mathbb{R}\).
\end{examsubparts}
\item
Let \((X,d)\) be a complete metric space. State and prove the contraction mapping theorem for~\(X\), which is also called the Banach fixed point theorem.
\end{examproblems}
\begin{examinstructions}
Answer the following with complete definitions or statements or short proofs.
\end{examinstructions}
\begin{examproblems}{5}
\item
Show that the axiom of choice is equivalent to the statement that for any indexed family \(\{A_\alpha\}_{\alpha\in J}\) of nonempty sets with \(J\ne\varnothing\), the cartesian product \(\prod_{\alpha\in J}A_\alpha\) is nonempty.
\item
Show that a connected normal space~\(X\) having more than one point is uncountable.
\item
Suppose that \(f:X\to Y\) is continuous and that \(A\subset X\) is connected. Show that \(f(A)\subset Y\) is connected.
\item
Suppose that \(f:X\to Y\) is continuous and that \(A\subset X\) is compact. Show that \(f(A)\) is compact.
\item
State the Extreme Value Theorem.
\item
Define what it means for a topological space to be regular.
\item
What does it mean for a topological space~\(X\) to be a Baire space?
\item
State the Cantor–Schröder–Bernstein theorem.
\item
Is every path connected space locally path connected?
\item
Define what it means for a topological space~\(X\) to be separable. Let~\(I\) denote the unit interval \([0,1]\). Is \(I^{\mathbb{Z}_+}\) separable? What about \(I^I\)?
\end{examproblems}
\end{document}
